% Sección a cargo de: Arteaga Alava Danela Dayana
\section{Ética y tratamiento de datos}\label{sec:etica}

Las entrevistas recogen la voz, las opiniones y, al cierre, el teléfono de los participantes. Son datos personales cuyo tratamiento regulan la Ley Orgánica de Protección de Datos Personales~\cite{lopdp2021} y su reglamento~\cite{reglamentolopdp2023}. El equipo aplica tres medidas: consentimiento informado antes de cada sesión, anonimización de lo que se publica y custodia restringida de los originales.

\subsection{Consentimiento informado}

El equipo elaboró el formulario PFC-ISR401-POSIA-A03 (versión 1.0 del 5 de noviembre de 2026), cuya versión pública está en el Anexo~\ref{anx:consentimiento}. El formulario explica el propósito, las actividades, los riesgos y los beneficios; reconoce los derechos de acceso, rectificación, eliminación y oposición de la LOPDP; y separa en nueve autorizaciones específicas lo que el participante acepta, de modo que puede, por ejemplo, participar sin autorizar la grabación. Incluye un desprendible para revocar el consentimiento y se firma en dos ejemplares, uno para el participante y otro para el expediente del proyecto.

En las tres entrevistas, realizadas el 19 de noviembre de 2026, el entrevistador leyó el guion de apertura del Anexo~\ref{anx:entrevista}, entregó el formulario y esperó la firma antes de comenzar con las preguntas. Los tres participantes autorizaron la grabación en audio.

Durante la revisión de las grabaciones se identificaron dos situaciones que el equipo corrige:

\begin{itemize}
\item \textbf{Grabación previa al consentimiento.} En la sesión ACTA-ELI-ADM-01 la grabadora quedó encendida antes de la lectura del consentimiento y captó una conversación previa ajena a la entrevista. Ese fragmento no se transcribe, no se analiza y se recorta del audio antes de su custodia. Desde ahora la grabación se inicia solo después de la firma.
\item \textbf{Dato no cubierto por las autorizaciones.} Al cierre de la sesión ACTA-ELI-CAJ-01 se pidió a la participante un correo de contacto, pero el formulario solo contempla el registro del número de teléfono (autorización e). El correo se elimina del registro y el guion de cierre se limita al teléfono, y solo si la autorización e) fue marcada.
\end{itemize}

\subsection{Anonimización}

\begin{enumerate}[label=\arabic*., leftmargin=*]
\item Cada sesión se identifica con un código (ACTA-ELI-CAJ-01, ACTA-ELI-CAJ-02, ACTA-ELI-ADM-01) y cada participante con su rol genérico (SH-01, SH-03, SH-04), como establece la sección~5 del formulario. Los nombres, cédulas, firmas y cargos específicos no aparecen en el documento ni en el repositorio.
\item En las transcripciones se reemplazan por \texttt{[establecimiento]}, \texttt{[compañera de trabajo]} o \texttt{[cliente]} el nombre del local, los nombres de terceros mencionados por los participantes y cualquier dato que permita identificarlos.
\item La identidad de la organización tampoco se publica: la versión pública del formulario omite su nombre y su dirección.
\item El documento cita hallazgos y paráfrasis, no fragmentos largos de las transcripciones. Una frase textual solo se cita si el participante marcó la autorización f).
\end{enumerate}

\subsection{Custodia de evidencias}

La custodia sigue lo comprometido en la sección~5 del formulario (Tabla~\ref{tab:custodia}).

\begin{table}[H]
\centering
\caption{Custodia de las evidencias de campo.}
\label{tab:custodia}
\small
\begin{tabularx}{\textwidth}{@{}L{3.6cm}L{3.6cm}X@{}}
\toprule
\textbf{Evidencia} & \textbf{Dónde se guarda} & \textbf{Regla} \\
\midrule
Formularios de consentimiento firmados y carta de aceptación & Archivo cifrado en la carpeta institucional restringida & Se escanean en PDF. La contraseña del archivo cifrado solo la conoce el docente responsable. \\
Grabaciones originales de audio & Mismo archivo cifrado & Nunca se suben al repositorio público. \\
Transcripciones anonimizadas & Carpeta institucional restringida & Se revisan con la regla de anonimización antes de compartirlas con el docente. \\
Registro de cada archivo & \texttt{evidencias/MANIFIESTO.md} en el repositorio & Nombre del archivo, fecha, duración y huella SHA-256, para comprobar que el archivo no cambió. \\
\bottomrule
\end{tabularx}
\end{table}

El material identificable se conserva como máximo 24 meses desde el cierre del periodo académico y luego se destruye. Solo se conservan los instrumentos en blanco y los resultados totalmente anonimizados. Si un participante entrega el desprendible de revocación, sus datos se eliminan siempre que aún no se hayan integrado en resultados agregados, y la eliminación se registra en el manifiesto.
